\NSPage{067}%
for this change of representation and preserves normalized Haar averages. Fix
\NSi{NS-F-P067-001}{(z,t)} with
\NSi{NS-F-P067-002}{0<q(z,t)<q_{\mathrm{big}}}. For a sufficiently small
neighborhood \NSi{NS-F-P067-003}{U} of this point, we set
\NSFormula{NS-F-P067-004}{display}{6.17}
\begin{equation}
  \begin{aligned}
    \mathcal D_\ell&=\supp_{(z,t)}\chi_\ell(q(z,t)),
      & \mathcal L(U)&=\{\ell:\mathcal D_\ell\cap U\ne\varnothing\},\\
    i_0&=\min_{\ell\in\mathcal L(U)}i(\ell),
      & H&=\pi_{i_0}(Y).
  \end{aligned}
  \tag{6.17}\label{eq:6.17}
\end{equation}
We choose \NSi{NS-F-P067-005}{U} so that its band indices differ by at most
four. Then \eqref{eq:6.14} gives the factorization
\NSFormula{NS-F-P067-006}{display}{6.18}
\begin{equation}
  \begin{aligned}
    \Delta_\ell&=i(\ell)-i_0\in\{0,\ldots,\Delta_{\max}\},
      & \pi_{i(\ell)}&=\pi_{\Delta_\ell}\circ\pi_{i_0},\\
    Y_i&=J_g^{\Delta_\ell}H,
      & \mathcal R_\gamma^H&=\pi_{\Delta_\ell}^{-1}(\mathcal R_\gamma),
      & \mathcal R_\gamma^{\mathrm{abs}}
        &=\pi_{i_0}^{-1}(\mathcal R_\gamma^H).
  \end{aligned}
  \tag{6.18}\label{eq:6.18}
\end{equation}
We define \NSi{NS-F-P067-007}{\mathcal R_\gamma^{H,+}} by the same preimage
formula with \NSi{NS-F-P067-008}{\mathcal R_\gamma^+}. Since
\NSi{NS-F-P067-009}{\pi_{i_0}} is surjective, the separation in
\eqref{eq:6.13} also holds for these common-torus rectangles. A sum of band
fields now has the representation
\NSFormula{NS-F-P067-010}{display}{none}
\[
  \sum_{\ell\in\mathcal L(U)}f_\ell(Y_{i(\ell)})
  =\sum_{\ell\in\mathcal L(U)}f_\ell(J_g^{\Delta_\ell}H).
\]

\begin{lemma}\label{lem:6.2}
For \NSi{NS-F-P067-011}{U,i_0,H} as in \eqref{eq:6.17}, each change
\NSi{NS-F-P067-012}{H\mapsto Y_i} has uniformly bounded derivative factors,
and \NSi{NS-F-P067-013}{\mathcal R_\gamma^H} consists of
\NSi{NS-F-P067-014}{14^{\Delta_\ell}\le 14^{\Delta_{\max}}} disjoint lifts
of \NSi{NS-F-P067-015}{\mathcal R_\gamma}. Haar averages on the absolute,
common, and band tori agree whenever the function is a pullback from the
respective torus. Radial integration at fixed \NSi{NS-F-P067-016}{(z,t)},
torus translation, averaging, and directional Fourier inversion on zero-mean
functions preserve common-torus descent and introduce no new dyadic band.
\end{lemma}

\begin{proof}
\textbf{Step 1: compare the covering maps and their averages.} We write
\NSi{NS-F-P067-017}{\Delta=\Delta_\ell}. The matrix norms of
\NSi{NS-F-P067-018}{J_g^\Delta} are bounded because
\NSi{NS-F-P067-019}{0\le\Delta\le\Delta_{\max}}. The preimage rectangles are
indexed by a finite quotient:
\NSFormula{NS-F-P067-020}{display}{none}
\[
  \begin{aligned}
  \mathcal R_\gamma^H
    &=\bigsqcup_{[k]\in\mathbb Z^2/J_g^\Delta\mathbb Z^2}
      \bigl\{J_g^{-\Delta}(c_\gamma+k+\xi v_r+\eta v_t)
      \pmod{\mathbb Z^2}:|\xi|,|\eta|<r_0\bigr\},\\
  \#(\mathbb Z^2/J_g^\Delta\mathbb Z^2)
    &=|\det J_g|^\Delta=14^\Delta.
  \end{aligned}
\]
Here \NSi{NS-F-P067-021}{c_\gamma} denotes the chosen representative in
\NSi{NS-F-P067-022}{\mathbb R^2}. Haar compatibility can be checked on
Fourier characters:
\NSFormula{NS-F-P067-023}{display}{6.19}
\begin{equation}
  \int_{\mathbb T^2} f(J_g^\Delta H)\,dH
  =\int_{\mathbb T^2} f(Y_i)\,dY_i.
  \tag{6.19}\label{eq:6.19}
\end{equation}
Indeed a character of frequency \NSi{NS-F-P067-024}{n} pulls back to
frequency \NSi{NS-F-P067-025}{(J_g^\Delta)^Tn}, which is zero exactly when
\NSi{NS-F-P067-026}{n=0}.

\textbf{Step 2: preserve the common representation under the stated
operations.} Radial integration fixes \NSi{NS-F-P067-027}{(z,t)}, and hence
\NSi{NS-F-P067-028}{q=q(z,t)}. It leaves every band cutoff unchanged even
though it can pass through many radial grid boxes. Torus translations commute
with deck translations; averaging and directional Fourier multipliers preserve
the lattice of pullback frequencies. These facts prove the preservation
claims.
\end{proof}

Fields are defined on the extended domain with the independent variable
\NSi{NS-F-P067-029}{Y}, using the original physical normalization; the common
torus is a local representation. On overlapping neighborhoods
\NSi{NS-F-P067-030}{U,V}, fix the same normalized band chart and write
\NSi{NS-F-P067-031}{f_U,f_V} for the two common-torus representatives. Their
compatibility is the identity
\NSFormula{NS-F-P067-032}{display}{6.20}
\begin{equation}
  f_U(R,Z,T,J_g^{i_0(U)}Y)=f_V(R,Z,T,J_g^{i_0(V)}Y)
  \tag{6.20}\label{eq:6.20}
\end{equation}
at points over \NSi{NS-F-P067-033}{U\cap V}. Both sides equal the original
field in that chart, so the identity also holds on the support closures where
the set of overlapping dyadic supports changes.

\NSPage{068}%
For a fixed band \NSi{NS-F-P068-001}{\ell\in\mathcal L(U)}, we write
\NSi{NS-F-P068-002}{\Delta=\Delta_\ell}. On its chart, a field descending to
\NSi{NS-F-P068-003}{H=J_g^{i_0}Y} has the operators in \eqref{eq:6.6} with
\NSi{NS-F-P068-004}{i} replaced by \NSi{NS-F-P068-005}{i_0}, while
\NSi{NS-F-P068-006}{Q} remains fixed. In particular,
\NSi{NS-F-P068-007}{c_{i_0}=T_g^{-\Delta}c_i} and
\NSi{NS-F-P068-008}{M_{i_0}=\Lambda_g^{-\Delta}M_i}. Since
\NSi{NS-F-P068-009}{\Delta} is bounded by \eqref{eq:6.18}, the change of torus
leaves the powers of \NSi{NS-F-P068-010}{\varepsilon} in the derivative bounds
unchanged.

Radial integration can collect more terms than pointwise multiplication. The
mesh gives the following bounds for sums over slow labels, with
\NSi{NS-F-P068-011}{S_\ell=\ell^2} when several bands are compared.

\begin{lemma}[Number of relevant labels]\label{lem:6.3}
There is a constant \NSi{NS-F-P068-012}{C}, independent of the band and point,
such that
\NSFormula{NS-F-P068-013}{display}{none}
\[
  \sup_{(r,z,t)}\#\{\gamma\in\Gamma:(r,z,t)\in K_\gamma\}\le C.
\]
For a fixed band \NSi{NS-F-P068-014}{\ell}, compact chart set
\NSi{NS-F-P068-015}{B\subset\mathbb R^3}, and bounded normalized radial
interval \NSi{NS-F-P068-016}{I_R},
\NSFormula{NS-F-P068-017}{display}{none}
\[
  \begin{aligned}
    \#\{(\ell,a,\sigma)\in\Gamma:B_{\ell,a}\cap B\ne\varnothing\}
      &\le C_BS_\ell^9,\\
    \sup_{Z,T}\#\{(\ell,a,\sigma)\in\Gamma:
      B_{\ell,a}\cap(I_R\times\{Z\}\times\{T\})\ne\varnothing\}
      &\le C_{I_R}S_\ell^3.
  \end{aligned}
\]
The same bounds hold for fixed-factor enlargements of the mesh boxes.
\end{lemma}

\begin{proof}
At a point, the dyadic partition and each one-dimensional grid partition have
bounded overlap; the two signs add a factor of two. On a bounded interval, a
mesh of size \NSi{NS-F-P068-018}{S_\ell^{-3}} has
\NSi{NS-F-P068-019}{O(S_\ell^3)} positions. A compact chart varies three
coordinates, giving \NSi{NS-F-P068-020}{O(S_\ell^9)} boxes; a radial line
varies only one, giving \NSi{NS-F-P068-021}{O(S_\ell^3)}. Fixed-factor
enlargement changes only the constants.
\end{proof}

During radial integration the possible bands are fixed:
\NSFormula{NS-F-P068-022}{display}{none}
\[
  \mathcal L(z,t)=\{\ell:(z,t)\in \mathcal D_\ell\},
  \qquad q=q(z,t)\text{ is independent of }r.
\]
The selected slow supports lie in a uniformly bounded normalized radial
interval. Thus Lemma~\ref{lem:6.3} bounds the number of labels met by the
integral by \NSi{NS-F-P068-023}{O(S_{\mathrm{ref}}^3)}, where
\NSi{NS-F-P068-024}{S_{\mathrm{ref}}=\ell_{\mathrm{ref}}^2} for any
\NSi{NS-F-P068-025}{\ell_{\mathrm{ref}}\in\mathcal L(U)}. All these
\NSi{NS-F-P068-026}{S_\ell} are comparable. Counting the common-torus lifts
adds at most the fixed factor \NSi{NS-F-P068-027}{14^{\Delta_{\max}}}. The sum
over these labels therefore permits polynomial growth in
\NSi{NS-F-P068-028}{S_*}, without changing the bands or their covering-index
differences.

To solve an amplitude equation from the beginning of a pulse to its current
coordinate \NSi{NS-F-P068-029}{v}, we need the earlier points on the same
lifted rectangle. The slow and transverse coordinates remain fixed along this
path. A source formed from several bands may be a function of
\NSi{NS-F-P068-030}{H} without being a function of a particular
\NSi{NS-F-P068-031}{Y_i}: it can take different values at the distinct
preimages of that band-torus point. We therefore describe the path separately
on each lifted rectangle. The linear functional
\NSi{NS-F-P068-032}{\lambda_t(w)=v_t\cdot w/(1+b_g^2)} satisfies
\NSi{NS-F-P068-033}{\lambda_t(v_r)=0},
\NSi{NS-F-P068-034}{\lambda_t(v_t)=1}. Within one lifted band rectangle, we
write
\NSi{NS-F-P068-035}{\eta_g(H)=\lambda_t(J_g^\Delta H-c_\gamma-k_{\mathrm{copy}})}.
For \NSi{NS-F-P068-036}{w\in[0,L_s]}, we define
\NSFormula{NS-F-P068-037}{display}{6.21}
\begin{equation}
  \begin{aligned}
    H_w&=H+T_g^{-\Delta}c_i\bigl(w-v(H)\bigr)v_t\\
       &=H_0+T_g^{-\Delta}c_iwv_t,
       & H_0&=H-T_g^{-\Delta}\bigl(\eta_g(H)+r_0\bigr)v_t.
  \end{aligned}
  \tag{6.21}\label{eq:6.21}
\end{equation}
This is the point with the same slow coordinates and band transverse
coordinate \NSi{NS-F-P068-038}{\xi_g}, and with pulse coordinate
\NSi{NS-F-P068-039}{w}. Thus a source
\NSi{NS-F-P068-040}{f(R,Z,T,H)} is evaluated along the path as
\NSi{NS-F-P068-041}{f(R,Z,T,H_w)}, without averaging over the other
preimages. Since \NSi{NS-F-P068-042}{\lambda_tJ_g^\Delta=T_g^\Delta\lambda_t},
at fixed \NSi{NS-F-P068-043}{w} one has
\NSFormula{NS-F-P068-044}{display}{6.22}
\begin{equation}
  D_HH_w=I-v_t\lambda_t,
  \qquad D_Hv=\frac{T_g^\Delta}{c_i}\lambda_t=O(S_*).
  \tag{6.22}\label{eq:6.22}
\end{equation}

\NSPage{069}%
Higher derivatives of these affine maps vanish. Equivalently, in the original
auxiliary coordinate \NSi{NS-F-P069-001}{Y}, the path is
\NSi{NS-F-P069-002}{Y'=Y+T_g^{-i}(\eta_g'-\eta_g)v_t}. A deck translation
preserving \NSi{NS-F-P069-003}{H} translates this path by the same amount and
reindexes the rectangle copy. A uniquely specified solution with zero initial
data therefore descends to the common torus. The precise differentiated
inverse estimate is proved in Proposition~\ref{prop:7.2}; \eqref{eq:6.22}
supplies its coordinate change without a new power of
\NSi{NS-F-P069-004}{\varepsilon}.

\subsection{Coefficient bounds and their algebra}\label{sec:6.4}

The residual estimates must account for both the size of a correction and its
vanishing at the boundary of its support. Let
\NSi{NS-F-P069-005}{f,g} denote angular mean coefficients and
\NSi{NS-F-P069-006}{a,b} denote labelled wave amplitudes. Their products and
normalized derivatives appear in the correction equations, for example
\NSFormula{NS-F-P069-007}{display}{none}
\[
  fg,\qquad fa,\qquad ab,\qquad D_rf,\qquad D_ra.
\]
We define three classes in Definitions~\ref{def:6.4} and~\ref{def:6.5} to
record the size of a mean coefficient, a wave amplitude, and a radial moment,
respectively:
\NSFormula{NS-F-P069-008}{display}{none}
\[
  f\in\mathsf M_\alpha,
  \qquad a_{\gamma,m}\in\mathsf W_\alpha,
  \qquad F(Z,T)\in\mathsf S_\alpha.
\]
The exponent \NSi{NS-F-P069-009}{\alpha} records a power of
\NSi{NS-F-P069-010}{\varepsilon}. We measure radial vanishing by the weight
\NSi{NS-F-P069-011}{\zeta} in \eqref{eq:6.23} and temporal decay by an
envelope \NSi{NS-F-P069-012}{P_v} satisfying
\NSi{NS-F-P069-013}{0<P_v\le1} on
\NSi{NS-F-P069-014}{0\le v\le L_s}. Its formula will be given in
\eqref{eq:7.12}. The estimates are imposed at every fixed derivative order so
that the shell extensions are smooth and the bounds survive the products and
derivatives in Proposition~\ref{prop:6.6}.

Recall the flat edge weight of Theorem~\ref{thm:4.6}. On the fixed active
interval \NSi{NS-F-P069-015}{X_a<X<X_b}, it may be written
\NSFormula{NS-F-P069-016}{display}{6.23}
\begin{equation}
  \zeta(X)=\exp\!\left(
    -\frac{a_a}{\log^2(X/X_a)}-\frac{a_b}{\log^2(X_b/X)}
  \right),
  \qquad
  \delta(X)=\min\{1,\log(X/X_a),\log(X_b/X)\},
  \tag{6.23}\label{eq:6.23}
\end{equation}
where \NSi{NS-F-P069-017}{a_a,a_b>0} are fixed; we extend
\NSi{NS-F-P069-018}{\zeta} by zero outside the interval. For every
\NSi{NS-F-P069-019}{b>0} and finite \NSi{NS-F-P069-020}{N},
\NSi{NS-F-P069-021}{\zeta^b\delta^{-N}} tends to zero faster than any power at
either boundary. Both \NSi{NS-F-P069-022}{\zeta} and
\NSi{NS-F-P069-023}{\delta} depend only on \NSi{NS-F-P069-024}{X}, so they are
constant along \eqref{eq:6.21}.

The rapidly varying phase and its amplitude have different derivative scales.
A \emph{coefficient derivative} acts on the amplitude with the oscillatory
exponential factored out. Fix a finite stage \NSi{NS-F-P069-025}{j}, and use
normalized chart representatives. For a coefficient
\NSi{NS-F-P069-026}{a}, we write
\NSFormula{NS-F-P069-027}{display}{none}
\[
  \partial^Ia
  =\partial_R^{I_1}\partial_Z^{I_2}\partial_T^{I_3}
   \partial_{H_1}^{I_4}\partial_{H_2}^{I_5}a,
  \qquad I\in\mathbb N_0^5.
\]
These derivatives hold the band, label, representative, and harmonic integer
fixed. The auxiliary coordinate derivatives measure each representative
separately and transform by the chain rule when \NSi{NS-F-P069-028}{i_0}
changes. In the definitions below, \NSi{NS-F-P069-029}{C_{j,I}>0} and
\NSi{NS-F-P069-030}{b_{j,I},d_{j,I}\ge0} may depend on
\NSi{NS-F-P069-031}{j,I} and the fixed profile data, and are uniform in band,
label, rectangle copy, and point.

\begin{definition}[Mean and moment coefficients]\label{def:6.4}
A smooth coefficient \NSi{NS-F-P069-032}{f=f(R,Z,T,H)} belongs to
\NSi{NS-F-P069-033}{\mathsf M_\alpha} if it descends to each applicable common
torus with representatives satisfying \eqref{eq:6.20}, extends smoothly by
zero outside the active radial shell, and on
\NSi{NS-F-P069-034}{X_a<X<X_b} satisfies
\NSFormula{NS-F-P069-035}{display}{6.24}
\begin{equation}
  \begin{gathered}
    \partial_\theta f=0,
    \qquad
    \forall I\in\mathbb N_0^5\ \exists C_{j,I},b_{j,I},d_{j,I}:\\
    |\partial^If|\le C_{j,I}\varepsilon^\alpha
      S_*^{b_{j,I}}\zeta\delta^{-d_{j,I}}.
  \end{gathered}
  \tag{6.24}\label{eq:6.24}
\end{equation}

\NSPage{070}%
It may depend on \NSi{NS-F-P070-001}{H}; it has no auxiliary rectangle
restriction or pulse-envelope bound. For a function
\NSi{NS-F-P070-002}{F=F(Z,T)}, define
\NSFormula{NS-F-P070-003}{display}{6.25}
\begin{equation}
  \begin{gathered}
    F\in\mathsf S_\alpha
    \quad\Longleftrightarrow\quad
    \forall I\in\mathbb N_0^2\ \exists C_{j,I},b_{j,I}:\\
    |\partial_{Z,T}^IF|\le C_{j,I}\varepsilon^\alpha S_*^{b_{j,I}}.
  \end{gathered}
  \tag{6.25}\label{eq:6.25}
\end{equation}
\end{definition}

Radial moment normalization. The moments to be estimated in
\NSi{NS-F-P070-004}{\mathsf S_\alpha} include the radial measure and weight in
their normalization: if
\NSi{NS-F-P070-005}{f_*=Q^{a_f}f_{\mathrm{phys}}}, then
\NSFormula{NS-F-P070-006}{display}{6.26}
\begin{equation}
  M_*(Z,T):=\int R^e\langle f_*\rangle_Y\,dR
  =Q^{a_f-(e+1)/2}\int r^e\langle f_{\mathrm{phys}}\rangle_Y\,dr.
  \tag{6.26}\label{eq:6.26}
\end{equation}
The auxiliary Haar average removes the dependence on
\NSi{NS-F-P070-007}{Y}; by \eqref{eq:6.19}, it can be taken on any applicable
common torus. The same scaling identity holds before averaging.

Phase, envelope, and support of a wave. For each label
\NSi{NS-F-P070-008}{\gamma}, we prescribe local data
\NSFormula{NS-F-P070-009}{display}{6.27}
\begin{equation}
  \begin{gathered}
    \Phi_\gamma=p_\gamma\theta+\varphi_\gamma(R,Z,T,H),
    \qquad k=\lceil\varepsilon^{-1/2}\rceil,
    \qquad kp_\gamma\in\mathbb Z\setminus\{0\},\\
    0<P_v\le1\quad(0\le v\le L_s),
    \qquad \mathscr H_j\subset\mathbb Z\setminus\{0\},
    \qquad |\mathscr H_j|<\infty.
  \end{gathered}
  \tag{6.27}\label{eq:6.27}
\end{equation}
Here \NSi{NS-F-P070-010}{p_\gamma} is fixed for each label,
\NSi{NS-F-P070-011}{\varphi_\gamma} is defined on its enlarged lifted
rectangles, and the harmonic set \NSi{NS-F-P070-012}{\mathscr H_j} is
independent of the band. The pulse construction specifies the phases and
envelopes in Equations~\eqref{eq:7.3} and~\eqref{eq:7.12}.

We write \NSi{NS-F-P070-013}{x_*=(R,Z,T)} and
\NSi{NS-F-P070-014}{K_\gamma^*=\mathcal C_\ell(K_\gamma)}. On each component
of \NSi{NS-F-P070-015}{\mathcal R_\gamma^H}, we use the coordinates
\NSi{NS-F-P070-016}{\xi_g,v} from \eqref{eq:6.11}. The domains and allowed
supports are
\NSFormula{NS-F-P070-017}{display}{6.28}
\begin{equation}
  \begin{aligned}
    \mathcal E_\gamma
      &=\{(x_*,H):H\in\mathcal R_\gamma^H\},\\
    \Omega_\gamma
      &=\{(x_*,H)\in\mathcal E_\gamma:x_*\in K_\gamma^*,
        \ \xi_g(H)\in\supp\chi_g,\ X_a\le X\le X_b\},\\
    \Omega_\gamma^{\mathrm{cut}}
      &=\Omega_\gamma\cap\{v(H)\in\supp\psi\}.
  \end{aligned}
  \tag{6.28}\label{eq:6.28}
\end{equation}
The slow coordinates range over the chart domain, and
\NSi{NS-F-P070-018}{X=QR^2/(2q)}. Replacing the closed rectangle in
\NSi{NS-F-P070-019}{\mathcal E_\gamma} by
\NSi{NS-F-P070-020}{\mathcal R_\gamma^{H,+}} gives the open enlarged domain
\NSi{NS-F-P070-021}{\mathcal E_\gamma^+}.

\begin{definition}[Labelled wave coefficients]\label{def:6.5}
For \NSi{NS-F-P070-022}{\gamma\in\Gamma} and
\NSi{NS-F-P070-023}{m\in\mathbb Z\setminus\{0\}}, a smooth amplitude
\NSi{NS-F-P070-024}{a_{\gamma,m}} on
\NSi{NS-F-P070-025}{\mathcal E_\gamma^+} belongs to
\NSi{NS-F-P070-026}{\mathsf W_\alpha} if its local formulas agree whenever
lifts represent the same common-torus point, satisfy \eqref{eq:6.20} on
overlapping neighborhoods, and
\NSFormula{NS-F-P070-027}{display}{none}
\[
  \partial_\theta a_{\gamma,m}=0,
  \qquad
  \supp(a_{\gamma,m}|_{\mathcal E_\gamma})\subset\Omega_\gamma.
\]
It extends smoothly by zero across the slow and transverse support boundaries
and both shell edges. On
\NSi{NS-F-P070-028}{\mathcal E_\gamma\cap\{X_a<X<X_b\}}, its bounds are
\NSFormula{NS-F-P070-029}{display}{6.29}
\begin{equation}
  \begin{gathered}
    \forall I\in\mathbb N_0^5\ \exists C_{j,I},b_{j,I},d_{j,I}:\\
    |\partial^Ia_{\gamma,m}|\le C_{j,I}\varepsilon^\alpha
      S_*^{b_{j,I}}\sqrt\zeta\,\delta^{-d_{j,I}}P_v,
    \qquad 0\le v\le L_s.
  \end{gathered}
  \tag{6.29}\label{eq:6.29}
\end{equation}
These are amplitude derivatives; \NSi{NS-F-P070-030}{P_v} is only a
pointwise weight. The continuation beyond
\NSi{NS-F-P070-031}{v=0,L_s} has no envelope bound or temporal zero-extension
condition.
\end{definition}

Cut-off fields and sums. An actual cut-off coefficient additionally has
\NSFormula{NS-F-P070-032}{display}{none}
\[
  \supp a_{\gamma,m}^{\mathrm{cut}}\subset\Omega_\gamma^{\mathrm{cut}}
  \qquad\text{and smooth zero extension in }v.
\]
No divisibility by \NSi{NS-F-P070-033}{\eta_\gamma,\chi_g,\psi} is assumed.

\NSPage{071}%
Let \NSi{NS-F-P071-001}{a_\nu} be a finite or locally finite family of
amplitudes, with labels \NSi{NS-F-P071-002}{\gamma_\nu\in\Gamma} and
harmonics \NSi{NS-F-P071-003}{m_\nu\in\mathscr H_j}, representing a real wave
\NSFormula{NS-F-P071-004}{display}{none}
\[
  w=\sum_\nu a_\nu e^{ik_{\gamma_\nu}m_\nu\Phi_{\gamma_\nu}}.
\]
Here \NSi{NS-F-P071-005}{k_\gamma} is the integer
\NSi{NS-F-P071-006}{k} in \eqref{eq:6.27} for the band of
\NSi{NS-F-P071-007}{\gamma}. We collect equal label--harmonic pairs before
testing membership in \NSi{NS-F-P071-008}{\mathsf W_\alpha}:
\NSFormula{NS-F-P071-009}{display}{none}
\[
  w=\sum_\gamma\sum_{m\in\mathscr H_j}
    A_{\gamma,m}e^{ik_\gamma m\Phi_\gamma},
  \qquad
  A_{\gamma,m}=\sum_{\nu:(\gamma_\nu,m_\nu)=(\gamma,m)}a_\nu.
\]
Then \NSi{NS-F-P071-010}{w\in\mathsf W_\alpha} means that every
\NSi{NS-F-P071-011}{A_{\gamma,m}} belongs to
\NSi{NS-F-P071-012}{\mathsf W_\alpha}, with the uniform bounds in
\eqref{eq:6.29}.

The bounds at every fixed derivative order justify smooth zero extension at
the shell edges: \NSi{NS-F-P071-013}{\zeta} and
\NSi{NS-F-P071-014}{\sqrt\zeta} absorb every allowed inverse power of
\NSi{NS-F-P071-015}{\delta}, including those introduced by differentiation.
We can now determine the size and support of the products that occur in the
momentum residual.

\begin{proposition}\label{prop:6.6}
Each class in Section~\ref{sec:6.4} is closed under finite sums at a fixed
exponent. For \NSi{NS-F-P071-016}{\alpha,\beta\in\mathbb R}, products satisfy
\NSFormula{NS-F-P071-017}{display}{6.30}
\begin{equation}
  \mathsf M_\alpha\mathsf M_\beta\subset\mathsf M_{\alpha+\beta},
  \qquad
  \mathsf M_\alpha\mathsf W_\beta\subset\mathsf W_{\alpha+\beta}.
  \tag{6.30}\label{eq:6.30}
\end{equation}
For \NSi{NS-F-P071-018}{a_{\gamma,m}\in\mathsf W_\alpha} and
\NSi{NS-F-P071-019}{b_{\gamma,m'}\in\mathsf W_\beta} with the same label,
\NSFormula{NS-F-P071-020}{display}{6.31}
\begin{equation}
  a_{\gamma,m}b_{\gamma,m'}\in
  \begin{cases}
    \mathsf W_{\alpha+\beta},&m+m'\ne0,
      \text{ with harmonic }m+m',\\
    \mathsf M_{\alpha+\beta},&m+m'=0,
      \text{ after temporal cutoff}.
  \end{cases}
  \tag{6.31}\label{eq:6.31}
\end{equation}
Before cutoff, a zero harmonic satisfies the mean bounds on its labelled
rectangle. Actual waves with distinct labels satisfy, for every pair of
derivative multi-indices \NSi{NS-F-P071-021}{I,J},
\NSFormula{NS-F-P071-022}{display}{none}
\[
  (\partial^Iw_\gamma)(\partial^Jw_{\gamma'})=0
  \qquad(\gamma\ne\gamma').
\]
For \NSi{NS-F-P071-023}{\mathcal C\in\{\mathsf M,\mathsf W\}}, every coefficient
derivative in a fixed common-torus representative preserves the
\NSi{NS-F-P071-024}{\varepsilon} exponent in its defining bounds. The
operators below also preserve \eqref{eq:6.20} and therefore act on the classes
by
\NSFormula{NS-F-P071-025}{display}{6.32}
\begin{equation}
  \begin{array}{c|c}
    \textit{operator}&\textit{mapping}\\ \hline
    \partial_{R,Z,T}&\mathcal C_\alpha\longrightarrow \mathcal C_\alpha\\
    D_r&\mathcal C_\alpha\longrightarrow \mathcal C_{\alpha-\kappa_s}\\
    D_z=\varepsilon\partial_Z&\mathcal C_\alpha\longrightarrow \mathcal C_{\alpha+1}\\
    Q^{1+h}N_{\mathrm{abs}}=c_{i_0}N_{i_0}
      &\mathcal C_\alpha\longrightarrow \mathcal C_\alpha\\
    -\varepsilon\partial_T&\mathcal C_\alpha\longrightarrow \mathcal C_{\alpha+1}.
  \end{array}
  \tag{6.32}\label{eq:6.32}
\end{equation}
These mappings concern amplitudes with the phase factored out. The local
coordinate derivative bounds, including those for
\NSi{NS-F-P071-026}{\partial_H}, retain the prescribed supports. The wave
rules are local to the labelled rectangles before cutoff, and hold for
smoothly extended fields on the common torus after cutoff. Products,
differentiation, and angular averaging preserve finite band-independent
harmonic sets at every fixed stage.
\end{proposition}

\begin{proof}
The triangle inequality gives closure under finite sums at a fixed exponent.

\textbf{Step 1: multiply the coefficient bounds.} Leibniz' formula bounds each
fixed derivative of a product using finitely many derivative bounds for its
factors. Powers of \NSi{NS-F-P071-027}{S_*} and
\NSi{NS-F-P071-028}{\delta^{-1}} add
\NSPage{072}%
and are allowed to increase. The weight bounds in Equations~\eqref{eq:6.30}
and~\eqref{eq:6.31} follow from \NSi{NS-F-P072-001}{0\le\zeta\le1},
\NSi{NS-F-P072-002}{0<P_v\le1}, and
\NSFormula{NS-F-P072-003}{display}{none}
\[
  \zeta^2\le\zeta,
  \qquad \zeta^{3/2}P_v\le\sqrt\zeta P_v,
  \qquad
  \zeta P_v^2\le
  \begin{cases}
    \zeta,&m+m'=0,\\
    \sqrt\zeta P_v,&m+m'\ne0.
  \end{cases}
\]
A mean coefficient may have larger radial or auxiliary support than a wave
coefficient. Their product is supported within the wave's support and retains
its envelope bound.

\textbf{Step 2: separate labels and select angular harmonics.} The support of
a derivative of a smoothly extended function is contained in its original
closed support. After the time cutoff has provided this smooth extension,
Lemma~\ref{lem:6.1} therefore eliminates all cross-label products of actual
fields. Before that cutoff, the coefficient algebra is used only on its
labelled band rectangle.

Recall that \NSi{NS-F-P072-004}{\langle\cdot\rangle_\theta} is the physical
angular average and \NSi{NS-F-P072-005}{\langle\cdot\rangle_Y} is normalized
Haar averaging in the auxiliary variable. By Lemma~\ref{lem:6.2}, the latter
can be computed in any applicable torus coordinates. Within one label the
product harmonic is \NSi{NS-F-P072-006}{m+m'}. Because
\NSi{NS-F-P072-007}{kp_\gamma} is a nonzero integer, angular averaging gives
\NSFormula{NS-F-P072-008}{display}{none}
\[
  \left\langle
    a_{\gamma,m}a_{\gamma,m'}e^{ik(m+m')\Phi_\gamma}
  \right\rangle_\theta
  =
  \begin{cases}
    a_{\gamma,m}a_{\gamma,m'},&m+m'=0,\\
    0,&m+m'\ne0.
  \end{cases}
\]
In particular,
\NSFormula{NS-F-P072-009}{display}{none}
\[
  \langle a_{\gamma,m}e^{ikm\Phi_\gamma}\rangle_\theta=0
  \quad(m\ne0),
  \qquad
  f^\circ=f-\langle f\rangle_Y
  \quad\text{for an angular mean }f.
\]
The second expression recalls the auxiliary zero-mean projection in
\eqref{eq:3.9}; it is different from both the physical angular mean and the
radial prefix average \NSi{NS-F-P072-010}{\mathcal A_X(f)} of
\eqref{eq:4.6}. A mean coefficient may therefore have a nonzero auxiliary
projection \NSi{NS-F-P072-011}{f^\circ}. Differentiation preserves each
harmonic integer and products take pairwise sums, so finitely many operations
preserve a finite harmonic set at a fixed stage.

\textbf{Step 3: apply the normalized derivative operators.} The coordinate
coefficients of \NSi{NS-F-P072-012}{D_r} and all their coefficient derivatives
are bounded by
\NSi{NS-F-P072-013}{C_I\varepsilon^{-\kappa_s}S_*^{C_I}} on the active
annulus. Together with \eqref{eq:6.6}, these bounds prove the exponent
estimates in \eqref{eq:6.32}; changing from a band to a common torus
contributes only the bounded matrices of Lemma~\ref{lem:6.2}. The local
\NSi{NS-F-P072-014}{\partial_H} bounds follow directly from \eqref{eq:6.24}
and \eqref{eq:6.29}.

In a fixed band chart, \NSi{NS-F-P072-015}{\partial_{R,Z,T}} preserves
\eqref{eq:6.20}, since the covering maps are independent of the slow
coordinates. The auxiliary terms
\NSi{NS-F-P072-016}{M_{i_0}\mathcal L_{i_0}} and
\NSi{NS-F-P072-017}{c_{i_0}N_{i_0}} represent the same absolute operators
\NSi{NS-F-P072-018}{Q^{d_r/2}\mathcal L_{\mathrm{abs}}} and
\NSi{NS-F-P072-019}{Q^{1+h}N_{\mathrm{abs}}} on every common torus. Thus the
remaining operators in \eqref{eq:6.32} also preserve compatibility. These
derivative rules concern amplitude coefficients; differentiating
\NSi{NS-F-P072-020}{e^{ikm\Phi_\gamma}} also differentiates the phase and
introduces its explicitly retained factor \NSi{NS-F-P072-021}{km}.

Differentiation, integration, translation, averaging, and Fourier multipliers
in \NSi{NS-F-P072-022}{(R,Z,T,H)} preserve independence of
\NSi{NS-F-P072-023}{\theta}.
\end{proof}

The wave support convention also fits the initial-value equations used to
construct pulses. The path \eqref{eq:6.21} fixes the same slow and transverse
variables. If a source vanishes there along the whole rectangle, its linear
solution with zero initial data is identically zero; after multiplication by
\NSi{NS-F-P072-024}{\psi}, it is supported in the original rectangle. Smooth
source extension and smooth dependence of the linear equation give smooth
extension at those boundaries. The mapping estimates at every fixed
derivative order are proved in Proposition~\ref{prop:7.2}. The
